\scp{\tsc{Wetar-Atauro}}{04-05}
The \tsc{Wetar-Atauro} node contains the languages spoken on
Wetar and Atauro islands,\footnote{
		There are also \ili{Kisar} speaking villages
		on Wetar island.}
as well as \ili{Galolen},
most speakers of which (c.
13,000) are located on the Timor mainland.
The \ili{Talur} dialect of \ili{Galolen} is spoken in three villages in south Wetar.
The other languages of Wetar are: \ili{Tugun},
\ili{Perai}, \ili{Aputai}, and Ili{\Q}uun.
Ili{\Q}uun is also spoken in the village of Ustutun on Lirang
Island which lies just to the southwest of Wetar.
This village has historically had an important role in the region,
as described by \citet[114]{Hinton2000}: “Ustutun is the most
influential village in the area
and the most progressive.
Because of the extreme tidal range,
the beach is used for boat repair by people from around the region”.
The languages of Atauro Island are \ili{Raklungu},
\ili{Hresuk}, \ili{Rasua}, and Dadu{\Q}a.
Dadu{\Q}a is also spoken on the northern coast of Timor,
which “is believed” to be due to a migration from Atauro Island
\citep[4]{Penn2006}, though the basis of this belief needs investigation.\footnote{
		The main source of data on Dadu{\Q}a is \citet{Penn2006} who describes the mainland variety.
		Speakers on the mainland are also bilingual in \ili{Galolen}
		and it may reflect some contact issues.}

\begin{table}\tcp{\tsc{Wetar-Atauro} *w}{04-18}
\begin{adjustbox}{width=\textwidth}
	\begin{threeparttable}\st{2pt}
	\begin{tabular}{llllllllll}\lsptoprule
*gloss		&	right								&	left								&	water						&	bee							&	root						&	eight							&	nine				&	snake				&	waist\\
PMP				&	*\tbr{w}anan				&	*\tbr{w}iRi					&	*\tbr{w}ahiyəR	&	*\tbr{w}ani			&	*\tbr{w}akaR		&	*\tbr{w}alu\su{‡}	&	*si\tbr{w}a	&	*sa\tbr{w}a	&	*ha\tbr{w}ak\\ \midrule
\ili{Tugun}			&	{\hr}malu\tbr{w}ana	&	{\hr}malu\tbr{w}iri	&	{\hr}eer				&	{\hr}ani				&	{\hr}akar				&	{\hr}fa-kau				&	{\hr}fa-sia	&	{\hr}sao		&	{\hr}oa-n\\
\ili{Perai}			&											&											&	{\hr}eer				&									&	{\hr}akar				&	{\hr}ha-kau				&	{\hr}sia		&	{\hr}sou		&	\\
\ili{Aputai}		&											&											&	{\hr}eer				&									&									&	{\hr}kau					&	{\hr}sia		&	{\hr}sao		&	\\
Ili{\Q}uun&	{\hr}\tbr{w}ana			&	{\hr}\tbr{w}iri			&	{\hr}eer				&	{\hr}ani				&	{\hr}akar				&	{\hr}kau					&	{\hr}ha-sia	&	{\hr}sao		&	\\
\ili{Hresuk}		&	{\hr}helana					&	{\hr}heliri					&	{\hr}eer				&	{\hr}ani				&	{\hr}aʔat				&	{\hr}he-ʔao				&	{\hr}he-see	&							&	\\
\ili{Raklungu}	&											&											&	\cg{(irai)}			&	{\hr}ani				&	{\hr}aʔa				&	{\hr}h-au					&	{\hr}ha-sia	&							&	\\
\ili{Rasua}			&	{\hr}palana					&	{\hr}paliri					&	{\hr}ee					&									&	{\hr}haʔa				&	{\hr}p-au					&	{\hr}pa-sia	&	{\hr}sau		&	\\
Dadu{\Q}a	&	{\hr}ma\tbr{w}ana		&	{\hr}ma\tbr{w}iik		&	{\hr}eeh				&	{\hr}ani				&									&	{\hr}wa-ʔao				&	{\hr}wa-sia	&							&	\\
\ili{Galolen}		&	{\hr}ma\tbr{w}ana		&	{\hr}ma\tbr{w}iik		&	{\hr}\tbr{w}ee	&	{\hr}\tbr{w}ani	&	{\hr}aʔar\su{†}	&	{\hr}i-hao				&	{\hr}i-sia	&	{\hr}sao		&	{\hr}oo\su{Δ}\\
		\lspbottomrule
	\end{tabular}
		\begin{tablenotes}\tnsize
			\item[†]	{\ili{Galolen} \it{aʔar} is attested in \citet{Hinton1989} from Hiay village on Wetar.
								Other known varieties of \ili{Galolen} have reflexes of *Ramut for ‘root’ (see \trf{04-23}).}
			\item[‡]	{Reflexes of ‘eight’ all have irregular initial \it{k} in Wetar
								languages and irregular *l > {\0} in all languages.
								Loss of medial *l in reflexes of *walu ‘eight’ also occurs in
								the \ili{Babar} islands and \tsc{\ili{Kisar}-Luang}.}
			\item[Δ]	{\ili{Galolen} \it{oo} means ‘body’.}
		\end{tablenotes}
	\end{threeparttable}
\end{adjustbox}
\end{table}

While there are sound changes shared between branches of \tsc{Wetar-Atauro}
(discussed below), there is only one (incomplete) sound change which encompass all
members of this group.
There is also one uniquely shared morphological innovation.
The shared sound change is *w > {\0},
though *w is retained unchanged in some words,
particularly word initially in \ili{Galolen}.
Examples of the reflexes of PMP *w are given in \trf{04-18}.
Note that loss of intervocalic *w is usually accompanied by assimilation of adjacent *a.

The morphological innovation found in \tsc{Wetar-Atauro} is **ka-
as a first person plural inclusive verbal subject agreement prefix.\footnote{
		Some languages of the \ili{Babar} Islands (\srf{04-06-05}) have first person
		plural inclusive prefixes with initial /k/,
		such as South\ili{east} \ili{Babar} \it{ke-} \citep[98]{Steinhauer2009}.
		These forms are a result of a regular *t > \it{k} sound change
		and are thus not cognate with the \tsc{Wetar-Atauro} prefix.}
 In \ili{Galolen} this prefix is only retained in the partially suppletive
inflection of \it{laʔa} ‘go’ with the first person inclusive
form \it{kaʔa}, as shown in \trf{04-19}.
The same form occurs in Dadu{\Q}a,
where it is also suppletive.
Reflexes of **ka- are retained productively in other
\tsc{Wetar-Atauro} languages (see \trf{04-20}).

\begin{table}\tcp{\tsc{Wetar-Atauro} *lakaw ‘go’ inflection}{04-19}
%\begin{adjustbox}{width=\textwidth}
	\st{4.5pt}\begin{tabular}{llllllll}\lsptoprule
	&	\tsc{1sg}	&	\tsc{2sg}	&	\tsc{3sg}	&	\tsc{1pl.incl}	&	\tsc{1pl.excl}	&	\tsc{2pl}	&	\tsc{3pl}\\	\midrule
\ili{Tugun} 	&	mu-laa	&	om-laa	&	ma-laa	&	\tbr{ka}-laa	&	am-laa	&	mar-laa	&	mar-laa\\
\tsc{(irrealis)}	&	&		&	&		&		&		&	\\
\ili{Hresuk}	&	laʔa	&	m-laʔa	&	n-laʔa	&	\tbr{k}-laʔa	&	m-laʔa	&	m-laʔa	&	k-laʔa\\
\ili{Raklungu}	&	laʔa	&	laʔa	&	laʔa	&	\tbr{k}-laʔa	&	k-laʔa	&	k-laʔa	&	k-laʔa\\
Dadu{\Q}a	&	laa	&	laa	&	laa	&	\tbr{k}aʔa	&	h-laa	&	h-laa	&	h-laa\\
\ili{Galolen}	&	laʔa	&	laʔa	&	laʔa	&	\tbr{k}aʔa	&	laʔa	&	laʔa	&	laʔa\\
		\lspbottomrule
	\end{tabular}
%\end{adjustbox}
\end{table}

Fuller verbal agreement paradigms for the \tsc{Wetar-Atauro}
languages are given in \trf{04-20}.
All languages have multiple paradigms which are partly determined
by the phonological shape of the stem.
\ili{Tugun} further has a distinction between agreement prefixes which
distinguish between realis and irrealis mood.\footnote{
Apart from the first plural person
inclusive slot, other parts of the \ili{Tugun} irrealis paradigm may have been formed
through fusion of agreement prefixes with \it{maa} ‘come’ from PMP *ma(R)i.
In Ili{\Q}uun cognate \it{maa} means ‘come’
and has an additional use described “as an auxiliary verb [which]
may simply state that something is going to happen in the (near)
future, but as a rule also expresses conditionality,
potentiality and (more weakly) intention of conditional exhortation.”
\citep[99]{DeJosselinDeJong1947}.}
Reflexes of \tsc{Wetar-Atauro}
first person inclusive **ka- are attested by \ili{Tugun} \tsc{1pl.incl.irr}
\it{ka-} and the \tsc{1pl.incl} prefix \it{k-} used by languages of Atauro
for most consonant initial stems.
In addition Ili{\Q}uun has a verbal prefix \it{ka-} which “always
refers to the first person plur[al] incl[usive] [{\ldots}] it
also almost invariably expresses adhortation or intention.”
\citep[98]{DeJosselinDeJong1947}.\footnote{
		Whether a cognate of first
		person plural inclusive **ka- is found in \ili{Aputai}
		and \ili{Perai} is unknown due to lack of data.
		In the absence of such data,
		we provisionally take the data from \ili{Tugun}
		and Ili{\Q}uun as representative of all Wetar languages.}
The parallels of this description with \ili{Tugun} \tsc{1pl.incl.irr} \it{ka-} should be obvious.

\begin{table}\tcp{\tsc{Wetar-Atauro} verb agreement}{04-20}
%\begin{adjustbox}{width=\textwidth}
	\begin{threeparttable}\st{5pt}
	\begin{tabular}{ll@{}llllllll}\lsptoprule
Language		&	env.					&				&	\tsc{1sg}&	\tsc{2sg}	&	\tsc{3sg}	&	\tsc{1pl.in}	&	\tsc{1pl.ex}	&	\tsc{2pl}	&	\tsc{3pl}	\\	\midrule
\multirow{2}{*}{\ili{Tugun} {$\left\{\hspace{-2mm}\begin{array}{c}\vp{|}\\\vp{|}\\\end{array}\hspace{-2mm}\right.$}}	&	{\#}C	&	\tsc{irr}	&	mu-	&	om-	&	ma-	&	ka- {\cgr}	&	am-	&	mar-	&	mar-	\\	
						&	{\#}C					&	\tsc{real}	
																		&	u-		&	o-		&	{\0}		&	it-							&	am-		&	mi-		&	ra-	\\
Ili{\Q}uun	&	{\#}C					&				&	{\0}	&	{\0}	&	{\0}		&	{\0}						&	{\0}	&	{\0}	&	{\0}	\\	
\ili{Hresuk}			&	{\#}C\su{†}		&				&	{\0}	&	m-		&	n-			&	k- {\cgr}				&	m-		&	m-		&	k-	\\	
\ili{Raklungu}		&	{\#}C\su{‡}		&				&	{\0}	&	{\0}	&	{\0}/n-	&	k- {\cgr}				&	k-		&	k-		&	k-	\\	
\ili{Rasua}				&	{\#}C\su{‡}		&				&	{\0}	&	{\0}	&	{\0}/n-	&	{\0}						&	{\0}	&	{\0}	&	{\0}	\\
Dadu{\Q}a		&	{\#}C\su{‡}		&				&	{\0}	&	{\0}	&	{\0}/n-	&	k-/{\0} {\cgr}	&	h-		&	h-		&	h-	\\	
\ili{Galolen}			&	{\#}C					&				&	{\0}	&	{\0}	&	{\0}		&	{\0}						&	{\0}	&	{\0}	&	{\0}	\\	\midrule
\multirow{2}{*}{\ili{Tugun} {$\left\{\hspace{-2mm}\begin{array}{c}\vp{|}\\\vp{|}\\\end{array}\hspace{-2mm}\right.$}}	&	{\#}V	&	\tsc{irr}	&	mu-	&	om-	&	man-	&	kat-	&	amʔ-	&	marr-	&	marr-	\\	
						&	{\#}V					&	\tsc{real}	
																		&	v-		&	m-		&	n-			&	itt-						&	amʔ-	&	mir-	&	r-	\\	
Ili{\Q}uun	&	{\#}V					&				&	{\0}	&	m-		&	n-			&	r-							&	r-		&	r-		&	r-	\\	
\ili{Hresuk}			&	{\#}V/h\su{Δ}	&				&	{\0}	&	m-		&	n-			&	r-							&	m-		&	m-		&	r-	\\	
\ili{Raklungu}		&	{\#}V/h\su{Δ}	&				&	m-		&	m-		&	n-			&	r-							&	r-		&	r-		&	r-	\\	
\ili{Rasua}				&	{\#}V					&				&	m-		&	m-		&	n-			&	r-							&	r-		&	r-		&	r-	\\	
Dadu{\Q}a		&	{\#}V					&				&	{\0}	&	m-		&	n-			&	r-							&	r-		&	r-		&	r-	\\	
\ili{Galolen}			&	{\#}V					&				&	{\0}	&	m-		&	n-			&	r-							&	r-		&	r-		&	r-	\\	
		\lspbottomrule
	\end{tabular}
		\begin{tablenotes}\tnsize
			\item[†]	{The \ili{Hresuk} {\#}C initial prefixes are only taken by stems with
								an initial sonorant /m n ŋ l/.
								Nasal prefixes are deleted or undergo metathesis with a following
								nasal resulting in the following changes: \it{n-+m → mn; n-/m-+ŋ
								→ ŋ; n-+n → n; m-+m → m}.}
			\item[‡]	{\ili{Raklungu}, \ili{Rasua},
								and Dadu{\Q}a replace initial /m/ with
								/n/ in the \tsc{3sg} form, e.g.
								Dadu{\Q}a \it{mate} ‘die’ \it{→ nii nate} ‘s/he dies’.
								In \ili{Raklungu} some, but not all,
								verbs with initial /r/ also replace this consonant with /n/, e.g.
								\ili{Raklungu} \it{riu} ‘bathe’ \it{→ nii niu} ‘s/he bathes’.}
			%\item[Δ]	{\ili{Raklungu} has an additional set of prefixes taken by verbs
								%with an initial consonant cluster
								%and a lexically specified set of verbs with an initial single consonant.
								%These prefixes are shaped \it{Ca-} where the initial consonant
								%corresponds to the prefixes used before vowel initial roots:
								%1/\tsc{2sg} \it{ma-}, \tsc{3sg} \it{na-}, PL \it{ra-}.}
			\item[Δ]	{Initial \it{h}
								and the prefix metathesise in \ili{Hresuk}
								and \ili{Raklungu} resulting in infixal allomorphs.}
		\end{tablenotes}
	\end{threeparttable}
%\end{adjustbox}
\end{table}

\clearpage
Apart from the two innovations which encompass all of \tsc{Wetar-Atauro},
there are sound changes defining sub-branches.
The languages of Atauro can be placed together on the basis of
shared *z > \it{r.} Additional evidence for an \tsc{Atauro} node
comes from the morphophonemic changes associated with \tsc{3sg}
agreement for stems beginning with /m/: metathesis in \ili{Hresuk} (e.g.
\it{n-+ maa → nii \tbf{mn}aa} ‘s/he comes’)
and /m/ → /n/ in Dadu{\Q}a and \ili{Raklungu} (e.g.
\ili{Raklungu} \it{n- + maa → nii \tbf{n}aa} ‘s/he comes’).
The Dadu{\Q}a and \ili{Raklungu} /m/ → /n/ replacement is probably
best explained as historic metathesis with subsequent cluster
reduction through deletion of the first consonant.
Within \tsc{Atauro}, \ili{Raklungu} and \ili{Hresuk} can be placed together
in a \tsc{\ili{Hresuk}-Raklungu} node on the basis of regular *b > \it{h}.
They further usually have *aw/wa > \it{a}.
\ili{Rasua} and Dadu{\Q}a are provisionally placed together in a \tsc{\ili{Rasua}-Dadu{\Q}a}
group on the basis of *b > \it{w}, though the scarcity of \ili{Rasua}
data means this node must remain tentative.

There is lexical evidence for placing the languages of Wetar
(apart from the \ili{Talur} dialect of \ili{Galolen}) in a single \tsc{Wetar}
node \citep{Hinton2000,Taber1993}.
The shared sound changes point most strongly to connecting \tsc{Wetar}
with \ili{Galolen}, but there are also changes linking \tsc{Wetar} to \tsc{Atauro};
that is there is evidence for a \tsc{Wetar-Galolen} node,
excluding \tsc{Atauro}, as well as a \tsc{Nuclear Wetar-Atauro} node, excluding \ili{Galolen}.
There are five sound changes in support of \tsc{Wetar-Galolen},
given in \qf{04-10}, and two in support of \tsc{Nuclear Wetar-Atauro}, given in \qf{04-11}.

{\wg\begin{xltabular}[l]{\linewidth}{>{\hspace{-\tabcolsep}}l<{\hspace{-\tabcolsep}}>{\hspace{-\tabcolsep}\enspace}
																				l<{\hspace{-\tabcolsep}}>{\hspace{-\tabcolsep}\space\space}
																				l<{\hspace{-\tabcolsep}}>{\hspace{-\tabcolsep}\space\space}
																				l
																				l>{\raggedright\arraybackslash}X}
\tbx{04-10}	&\mc{5}{l}{\tsc{Wetar-Galolen} innovations}\\
				&a.	&*z			&>&\it{s}	&(\tsc{Atauro} *z > \it{r})\\
				&b.	&*ŋ/*n	&>&\it{n}	&(Dadu{\Q}a, *ŋ > \it{n}; \ili{Rasua}, \ili{Hresuk}, \ili{Raklungu} *ŋ = \it{ŋ})\\
				&c.	&*b			&>&**f		&(subsequent **f > \it{h} in Ili{\Q}uun and \ili{Galolen}; Dadu{\Q}a, \ili{Rasua} *b > w; \ili{Hresuk}, \ili{Raklungu} *b > h)\\
				&d.	&*-ay		&>&\it{e}	&(Dadu{\Q}a, \ili{Hresuk} *ay > \it{e}; \ili{Raklungu}, \ili{Rasua} *ay > \it{i})\\
				&e.	&*-aw		&>&\it{o}	&(Dadu{\Q}a, \ili{Rasua}, \ili{Hresuk} *aw > \it{a, o}; \ili{Raklungu} *aw > \it{a})\\ \midrule\hhline{~}
%\end{xltabular}\addtocounter{table}{-1}}
%
%{\wg\tblr{>{\hspace{-\tabcolsep}}l<{\hspace{-\tabcolsep}}>{\hspace{-\tabcolsep}\enspace}
																				%l<{\hspace{-\tabcolsep}}>{\hspace{-\tabcolsep}\space\space}
																				%l<{\hspace{-\tabcolsep}}>{\hspace{-\tabcolsep}\space\space}
																				%l
																				%l>{\raggedright\arraybackslash}X}
\tbx{04-11}	&\mc{5}{l}{\tsc{Nuclear Wetar-Atauro} innovations}\\
				&a.	&*j			&>&{\0}, \it{r}					&(\ili{Galolen} *j > \it{l, r})\\
				&b.	&*R/*d	&>&\it{r}								&(\ili{Galolen} *R > {\0})\\
				&c.	&*t			&>&\it{k, t} /{\gap\#}	&(\ili{Galolen} *t = \it{t} /{\gap}{\#})\\
%\end{tabular}\mbox{}\\}
\end{xltabular}\addtocounter{table}{-1}}

Of the changes defining \tsc{Wetar-Galolen},
all except *z > \it{s} are regionally common
and most are also found in Dadu{\Q}a.
The change *z > \it{s} occurs in the \tsc{\ili{Teun}-\ili{Nila}-Serua} of
\tsc{Southwest Maluku} (\srf{04-06-04}), as well as in some languages of \tsc{\ili{Central} Timor} (not a member
of \tsc{Timor-Babar}, see \crf{ch:CT}).
The change *ŋ > \it{n} has occurred in almost all \tsc{Timor-Babar}
languages to various extents
and probably spread by diffusion throughout this region.
Examples of the reflexes of *z,
*ŋ, and *b in \tsc{Wetar-Atauro} are given in \trf{04-21}.

\begin{table}
	\centering\tcp{\tsc{Wetar-Atauro} *z, *ŋ, and *b}{04-21}
	\st{3.2pt}\begin{tabular}{llll|lll|ll}\lsptoprule
*gloss	&	rain	&	path	&	far	&	name	&	wind	&	ten	&	stone	&	moon\\
PMP	&	*qu\tbr{z}an	&	*\tbr{z}alan	&	*\tbr{z}auq	&	*\tbr{ŋ}ajan	&	*ha\tbr{ŋ}in	&	*sa\tbr{ŋ}apuluq	&	*\tbr{b}atu	&	*\tbr{b}ulan\\	\midrule
\ili{Raklungu}	&	{\hr}u\tbr{r}a	&	{\hr}\tbr{r}ala	&	{\hr}\tbr{r}oo	&	{\hr}\tbr{ŋ}ai	&	{\hr}a\tbr{ŋ}i	&	{\hr}sa\tbr{ŋ}ulu	&	{\hr}\tbr{h}atu	&	{\hr}\tbr{h}ula\\
\ili{Hresuk}	&	{\hr}u\tbr{r}an	&	{\hr}\tbr{r}alan	&	{\hr}\tbr{r}oo	&	{\hr}\tbr{ŋ}ae	&	{\hr}a\tbr{ŋ}in	&	{\hr}se\tbr{ŋ}ulu	&	{\hr}\tbr{h}atu	&	{\hr}\tbr{h}ulan\\
\ili{Rasua}	&	{\hr}u\tbr{r}a	&	{\hr}\tbr{d}ala	&	{\hr}ra\tbr{r}oo	&	{\hr}\tbr{ŋ}ai	&	{\hr}a\tbr{ŋ}i	&	{\hr}sa\tbr{n}ulu	&	{\hr}\tbr{w}atu	&	{\hp{*h}}ula\\
Dadu{\Q}a	&	{\hr}u\tbr{r}a	&	{\hr}\tbr{r}ala	&	{\hr}\tbr{r}oo	&	{\hr}\tbr{n}ae	&	{\hr}a\tbr{n}i	&	{\hr}sa\tbr{n}ulu	&	{\hr}\tbr{w}atu	&	{\hr}\tbr{w}ula\\
Ili{\Q}uun	&	{\hr}u\tbr{s}an	&	{\hr}\tbr{s}alan	&	{\hr}\tbr{s}oo	&	{\hr}\tbr{n}ean	&	{\hr}a\tbr{n}in	&	{\hr}sa\tbr{n}ulu	&	{\hr}\tbr{h}atu	&	{\hr}\tbr{h}ulan\\
\ili{Perai}	&	{\hr}u\tbr{s}a	&	{\hr}\tbr{s}ala	&	{\hr}\tbr{s}oo	&	{\hr}\tbr{n}eaŋ	&	{\hr}a\tbr{n}in	&	{\hr}sa\tbr{n}ulu	&	{\hr}\tbr{f}atu	&	{\hr}\tbr{f}ulaŋ\\
\ili{Aputai}	&	{\hr}u\tbr{s}a	&	{\hr}\tbr{s}ala	&	{\hr}\tbr{s}oo	&	{\hr}\tbr{n}eaŋ	&	{\hr}a\tbr{n}iŋ	&	{\hr}sa\tbr{n}uluŋ	&	{\hr}\tbr{f}atu	&	{\hr}\tbr{f}ulan\\
\ili{Tugun}	&	{\hr}u\tbr{s}a	&	{\hr}\tbr{s}ala	&	{\hr}\tbr{s}oo	&	{\hr}\tbr{n}ean	&	{\hr}a\tbr{n}in	&	{\hr}sa\tbr{n}ulu	&	{\hr}\tbr{f}atu	&	{\hr}\tbr{f}ulan\\
\ili{Galolen}	&	{\hr}u\tbr{s}an	&	{\hr}\tbr{s}alan	&	{\hr}\tbr{s}oo	&	{\hr}\tbr{n}aran	&	{\hr}a\tbr{n}in	&	{\hr}sa\tbr{n}ulu	&	{\hr}\tbr{h}atu	&	{\hr}\tbr{h}ulan\\
		\lspbottomrule
	\end{tabular}
\end{table}

Of the sound changes shared between \tsc{Wetar}
and \tsc{Atauro} the split of *j > {\0}\tsc{+fr},
\it{r} has the same outcome in each language,
with the exception of reflexes of *huaji which exceptionally
has *j > \it{l} in \tsc{Wetar}.
\ili{Galolen} has a split of *j > \it{r,
l} which does not align in a systematic way with the \tsc{Nuclear
Wetar-Atauro} split.\footnote{
		Waima{\Q}a and \ili{Naueti} have \it{ʔlara} ‘sun’
		and \ili{Midiki} has \it{lara} ‘sun’.}
Examples of the reflexes of
*j and *w in \tsc{Wetar-Atauro} are given in \trf{04-22}.
Two additional examples of \ili{Galolen} *j are *hapəjis > **ma-pəjis
> \it{beli} ‘sick, pain’ and *bujəq > \it{na-hulan} ‘foam (verb)’.

\begin{table}\tcp{\tsc{Wetar-Atauro} *j}{04-22}
\begin{adjustbox}{width=\textwidth}
	\begin{threeparttable}\st{2pt}
	\begin{tabular}{llllllll}\lsptoprule
local gl.		&	sun								&	sick, pain\su{†}	&	name						&	dry\su{‡}						&	nose								&	how m.?							&	ySib\\
PMP					&	*qalə\tbr{j}aw		&	*qapə\tbr{j}u			&	*ŋa\tbr{j}an		&	*ma\tbr{j}a					&	*(ŋ)i\tbr{j}uŋ\su{Δ}&	*pi\tbr{j}a					&	*hua\tbr{j}i\su{ǁ}\\	\midrule
\ili{Raklungu}		&	{\hr}lalea				&	{\hqa}peu					&	{\hr}ŋai				&	\cg{(rose)}					&	{\hr}hŋi\tbr{r}u		&	{\hr}ha-i\tbr{r}aʔ	&	{\hr}ai, aik\\
\ili{Hresuk}			&	\hp{*la}lea				&	{\hqa}peu					&	{\hr}ŋae				&	\cg{(rosen)}				&	{\hr}ŋi\tbr{r}uk		&	{\hr}heʔ-i\tbr{r}a	&	{\hr}aik\\
\ili{Rasua}				&	\hp{*la}lea				&	{\hqa}peu					&	{\hr}ŋai				&	\cg{(rohe)}					&	{\hr}ŋi\tbr{r}u			&	{\hr}w-i\tbr{r}a		&	\\
Dadu{\Q}a		&	\hp{*la}lea				&	{\hqa}beu					&	{\hr}nae				&	\cg{(nagenah)}			&	{\hr}ni\tbr{r}u			&	{\hr}waʔ-i\tbr{r}a	&	{\hr}a\tbr{l}ik\\
Ili{\Q}uun	&	\hp{*la}leo				&	{\hqa}peu					&	{\hr}nean				&	{\hr}mea						&	\cg{(kunai)}				&	{\hr}haʔ-i\tbr{r}a	&	{\hr}a\tbr{l}i, a\tbr{l}in\\
\ili{Perai}				&	\hp{*la}leo				&										&	{\hr}neaŋ				&	\cg{(geraŋ)}				&	{\hr}nunu\tbr{r}		&	{\hr}faʔ-i\tbr{r}a	&	{\hr}a\tbr{l}i\\
\ili{Aputai}			&	\hp{*la}leo				&										&	{\hr}neaŋ				&	\cg{(geraŋ)}				&	{\hr}nunu\tbr{r}		&	{\hr}faʔ-i\tbr{r}a	&	{\hr}a\tbr{l}i\\
\ili{Tugun}				&	\hp{*la}leo				&	{\hqa}peu					&	{\hr}nean				&	\cg{(geran)}				&	\cg{(kunai)}				&	{\hr}faʔ-i\tbr{r}a	&	{\hr}a\tbr{l}i\\
\ili{Galolen}			&	\hp{*la}le\tbr{l}o&	{\hqa}e\tbr{l}u-n	&	{\hr}na\tbr{r}an&	{\hr}ra-ma\tbr{r}an	&	{\hr}inu\tbr{r}			&	{\hr}i-hi\tbr{l}a		&	{\hr}a\tbr{l}ir\\
		\lspbottomrule
	\end{tabular}
		\begin{tablenotes}\tnsize
			\item[†]	{\ili{Galolen} \it{elu-n} retains the meaning ‘gall,
								bile’ which is also reconstructed for PMP *qapəju.
								Other languages have reflexes of intermediate **ma-pəju.}
			\item[‡]	{Ili{\Q}uun \it{mea} is glossed as ‘dried up,
								coagulated’ \citep[128]{DeJosselinDeJong1947}
								and \ili{Galolen} \it{ra-maran} ‘make something dry’ (Indonesian \it{mengering}, \ili{Tetun} \it{hamaran}).}
			\item[Δ]	{\ili{Galolen} \it{inur} ‘nose’ shows consonant metathesis from
								intermediate **irun and/or could be a borrowing from \ili{Tetun} \it{inur} ‘nose’.
								\citet{BlustTrussel2020} also reconstruct PMP *ujuŋ ‘nose’
								and initial *u in this word may explain the initial vowel of
								\ili{Perai} and \ili{Aputai} \it{nunur} ‘nose’ These words may also show metathesis
								of the medial and final consonants.}
			\item[{ǁ}]	{\ili{Raklungu} \it{ai} ‘same sex sibling (unmarked for relative
								age.)’, Ili{\Q}uun \it{ali, alin} ‘same sex younger sibling’.}
		\end{tablenotes}
	\end{threeparttable}
\end{adjustbox}
\end{table}

The change of word-final *t > \it{k,
t} /{\gap}{\#} which is shared by \tsc{Wetar}
and \tsc{Atauro} appears to have spread by diffusion.\footnote{
		\citet[382]{Blust2004}
		considers the possibility that *t > \it{k} /{\gap}{\#} spread
		by diffusion or arose once in the immediate ancestors of the
		Wetar languages and tentatively takes the latter position.}
This change is not
complete and counter-examples can be found.
Examples of *t > \it{k,
t} /{\gap}{\#} are given in \trf{04-23}.
Given that this change is complete in Ili{\Q}uun,
combined with the socially prominent role of Ili{\Q}uun on Lirang
island \citep[114]{Hinton2000}, it is likely that Ili{\Q}uun was the centre of diffusion for
final *-t > \it{k}.

\begin{table}\tcp{\tsc{Wetar-Atauro} *t > \it{k, t} /{\gap\#}}{04-23}
%\begin{adjustbox}{width=\textwidth}
	\begin{threeparttable}\st{2pt}
	\begin{tabular}{llllllll}\lsptoprule
local gl.		&	skin						&	root									&	heavy								&	leech							&	vein									&	west							&	four\\
PMP					&	*kuli\tbb{t}		&	*Ramu\tbb{t}					&	*(ma)-bəRəqa\tbb{t}	&	*ma\tbb{t}ək\su{†}&	*uRa\tbb{t}						&	*habaRa\tbb{t}		&	*əpa\tbb{t}\\	\midrule
\ili{Raklungu}		&	{\hr}uli\tbr{k}	&	\cg{(aʔa)}						&											&										&												&	{\hr}hara\tbr{k}	&	{\hr}h-aaʔ\\
\ili{Hresuk}			&	{\hr}uli\tbr{k}	&	\cg{(aʔat)}						&	{\hr}pera\tbb{t}		&										&	{\hr}ura\tbr{k}				&										&	{\hr}haʔ-aa\tbb{t}\\
%\ili{Rasua}				&	{\hr}uli				&	\cg{(haʔa)}						&	{\hr}pera						&										&												&										&	{\hr}w-aa\\
Dadu{\Q}a		&	{\hr}uli\tbr{k}	&	{\hr}ramu\tbr{k}			&	{\hr}bera\tbr{k}		&	\cg{(urdua)}			&												&										&	{\hr}waa\tbr{k}\\
Ili{\Q}uun	&	{\hr}uli\tbr{k}	&	\cg{(akar)}						&	{\hr}hera\tbr{k}		&	{\hr}maka\tbr{k}	&	{\hr}ura\tbr{k}				&	{\hr}hara\tbr{k}	&	{\hr}aa\tbr{k}\\
\ili{Perai}				&	{\hr}eli\tbr{k}	&	\cg{(akar)}						&											&	{\hr}maka\tbr{k}	&												&										&	{\hr}h-aa\tbr{k}\\
\ili{Aputai}			&	{\hr}eli\tbr{k}	&	{\hr}ramu\tbr{k}			&											&	{\hr}maka\tbr{k}	&												&										&	{\hr}aa\tbr{k}\\
\ili{Tugun}				&	{\hr}uli\tbr{k}	&	\cg{(akar)}						&	{\hr}fera\tbr{k}		&	{\hr}maka\tbb{t}	&	{\hr}ua\tbb{t}\su{‡}	&	{\hr}fara\tbb{t}	&	{\hr}f-aa\tbb{t}\\
\ili{Galolen}, 		&	{\hr}uli\tbb{t}	&	{\hr}amu\tbr{k}, 			&											&	{\hr}maka\tbb{t}	&												&										&	{\hr}in-haa\tbb{t}\\
\ili{Talur}				&									&	{\hr}amu\tbb{t}\su{Δ}	&											&										&												&										&	\\
\ili{Galolen}			&	{\hr}uli\tbb{t}	&	{\hr}amu\tbb{t}				&											&	{\hr}ma\tbb{t}ak	&												&										&	{\hr}i-haa\tbb{t}\\
		\lspbottomrule
	\end{tabular}
		\begin{tablenotes}\tnsize
			\item[†]	{Reflexes in Wetar have consonant metathesis from PMP *matək
								and are thus from intermediate **makət.}
			\item[‡]	{\ili{Tugun} \it{uat} ‘vein’ has irregular *R > {\0} (expect *R > \it{r}).}
			\item[Δ]	{\ili{Galolen} from Ilputih village has \it{amuk} ‘root’,
								\ili{Galolen} from Ilwaki village has \it{amut} ‘root’.}
		\end{tablenotes}
	\end{threeparttable}
%\end{adjustbox}
\end{table}

Another \tsc{Nuclear Wetar-Atauro} change is the merger of *R/*d > \it{r}.
Examples of PMP *R in \tsc{Wetar-Atauro} are given in \trf{04-24}.
Two words have *R > {\0} in all known \tsc{Wetar-Atauro} reflexes:
*laRiw ‘run’ > \ili{Rasua},
Dadu{\Q}a, \ili{Galolen} \it{lai}, \ili{Raklungu} \it{kaklai},
\ili{Hresuk} \it{keklai}, and *diRus ‘bathe’ > \ili{Hresuk} \it{k-rius,}
\ili{Raklungu} \it{k-riu}, \ili{Rasua} \it{riu}, Dadu{\Q}a \it{riuh}, \ili{Galolen} \it{rius}.
One word has *R = \it{r} in all reflexes: *Ribu ‘thousand’ >
\ili{Rasua}, \ili{Hresuk}, \ili{Galolen}, Ili{\Q}uun, \ili{Aputai},
\ili{Tugun} \it{rihun}, Dadu{\Q}a \it{rihu},
\ili{Perai} \it{riuŋ}.

\begin{table}\tcp{\tsc{Wetar-Atauro} *R}{04-24}
\begin{adjustbox}{width=\textwidth}
	\st{2pt}\begin{tabular}{llllllll}\lsptoprule
local gl.	&	red	&	new	&	thorn, bone	&	blood	&	soil	&	descend	&	house\\
PMP	&	*ma-i\tbr{R}aq	&	*baqə\tbr{R}u	&	*du\tbr{R}i	&	*da\tbr{R}aq	&	*da\tbr{R}əq	&	*tu\tbr{R}un	&	*\tbr{R}umaq\\ \midrule
\ili{Raklungu}	&	{\hr}ma{\tl}me\tbr{r}a	&	{\hr}ha{\tl}he\tbr{r}u	&	{\hr}ru\tbr{r}i	&	{\hr}ra\tbr{r}a	&	{\hr}ra\tbr{r}i	&	{\hr}tu\tbr{r}u	&	{\hr}\tbr{r}uma\\
\ili{Hresuk}	&	{\hr}me{\tl}me\tbr{r}a	&	{\hr}he{\tl}he\tbr{r}u	&	{\hr}ru\tbr{r}i-k	&	{\hr}ra\tbr{r}a-k	&	{\hr}ra\tbr{r}e	&	{\hr}tu\tbr{r}un	&	{\hr}\tbr{r}uma\\
\ili{Rasua}	&	{\hr}ma{\tl}me\tbr{r}a	&	{\hr}wa{\tl}we\tbr{r}u	&	{\hr}ru\tbr{r}i	&	{\hr}ra\tbr{r}a	&	{\hr}ra\tbr{r}i ihi	&		&	{\hr}\tbr{r}uma\\
Dadu{\Q}a	&	{\hr}ma{\tl}me\tbr{r}a	&	{\hr}wa{\tl}we\tbr{r}u	&	{\hr}ru\tbr{r}i	&	{\hr}ra\tbr{r}a	&	{\hr}ra\tbr{r}e	&	{\hr}tu\tbr{r}u	&	{\hr}\tbr{r}uma\\
Ili{\Q}uun	&	{\hr}me\tbr{r}a	&	{\hr}ha{\tl}he\tbr{r}un	&	{\hr}ru\tbr{r}i-n	&	{\hr}ra\tbr{r}a	&	{\hr}ra\tbr{r}e	&	{\hr}tu\tbr{r}u\cg{(n)}	&	\cg{(raha)}\\
\ili{Perai}	&	{\hr}me\tbr{r}a	&	{\hr}fe\tbr{r}u	&	{\hr}ru\tbr{r}i-ŋ	&	{\hr}ra\tbr{r}a-ŋ	&	{\hr}re\tbr{r}e	&		&	\cg{(raha)}\\
\ili{Aputai}	&	{\hr}me\tbr{r}a	&	{\hr}fe\tbr{r}u	&	{\hr}ru\tbr{r}i-ŋ	&	{\hr}ra\tbr{r}a-ŋ	&	{\hr}ra\tbr{r}e	&		&	\cg{(raha)}\\
\ili{Tugun}	&	{\hr}me\tbr{r}a	&	{\hr}fe\tbr{r}u	&	{\hr}ru\tbr{r}i-n	&	{\hr}ra\tbr{r}a	&	{\hr}ra\tbr{r}e, ra\tbr{r}i	&	{\hr}tu\tbr{r}u	&	\cg{(raha)}\\
\ili{Galolen}	&	{\hr}mea	&	{\hr}heu	&	{\hr}rui-n	&	{\hr}raa-n	&	{\hr}rea	&	{\hr}tuun	&	{\hr}uma\\
		\lspbottomrule
	\end{tabular}
\end{adjustbox}
\end{table}

To summarise, there are four sound changes in \tsc{Wetar-Galolen}
which do not appear to be the result of diffusion: *z > \it{s};
*b > *f; *ay > \it{e}; *aw > \it{o},
and two sound changes in the other \tsc{Wetar-Atauro} languages
which may not be diffusional: *j > {\0}\tsc{+fr},
\it{r} and *R/*d > \it{r}.

Given this, we tentatively propose that the first split in the
\tsc{Wetar-Atauro} family was between \tsc{Atauro} and \tsc{Wetar-Galolen}.
After this split, the changes defining \tsc{Wetar-Galolen} then developed.
\ili{Galolen} then split from \tsc{Wetar}
and the changes in \tsc{Nuclear Wetar-Atauro} developed
and spread between the languages of \tsc{Wetar} and \tsc{Atauro}.
The lexical similarity of the languages of Wetar either indicates
subsequent co-development or that earlier diversity was erased
by later dialect/language levelling.
Throughout this period \ili{Galolen} maintained links with Wetar,
as attested by the \ili{Talur} dialect of \ili{Galolen} spoken in the south
of the island with on-going contact with the mainland
(João Cristo Rei pers. comm. 2011).
Whether \ili{Talur} is a remnant left after migration to Timor or a
result of a later back-migration is unknown.

This scenario is not to deny the reality that \tsc{Wetar} subgroups
with both \tsc{Atauro} and \ili{Galolen} and has undergone a period of common development with both groups.
Instead, we are proposing that the period of common development
shared between \tsc{Wetar-Galolen} occurred before the period
shared between \tsc{Nuclear Wetar-Atauro}.
A simplified wave diagram of the \tsc{Wetar-Atauro} languages
is given in \frf{04-03}.

\begin{figure}\centering
	\includegraphics[width=0.5\columnwidth]{Wetar-Atauro-wave.pdf}
	\fcp{Simplified \tsc{Wetar-Atauro} wave diagram}{04-03}
\end{figure}
